\chapter*{Apéndices}
\addcontentsline{toc}{chapter}{Apéndices}
\label{inicio_anexos}
Este capítulo reúne los anexos que el Manual cita en sus capítulos 1 a 5. Cada anexo se identifica como Anexo 1 a Anexo 12. Los formatos tienen el mismo código que los del Manual de Operaciones (F-001 a F-017) y del Manual de Control de Mantenimiento (F-018 a F-022).

\renewcommand{\arraystretch}{1.25}

% =====================================================================
\section*{Anexo 1 --- Objetivos de Seguridad Operacional y su Relación con los SPI}
\addcontentsline{toc}{section}{Anexo 1 --- Objetivos de Seguridad Operacional y SPI}
Cada objetivo específico del capítulo 2 se mide con uno o más indicadores (SPI, ver Anexo 9). El Gerente SMS revisa esta relación una vez al año, al revisar los objetivos (ver el capítulo 2), y cada vez que cambie el alcance de la operación.

\begin{longtable}{|>{\RaggedRight\arraybackslash}p{4cm}|>{\RaggedRight\arraybackslash}p{5.2cm}|>{\RaggedRight\arraybackslash}p{6.6cm}|}
\hline
\rowcolor[gray]{0.9}\textbf{Objetivo específico} & \textbf{SPI asociados (Anexo 9)} & \textbf{Por qué se mide así} \\ \hline
\endhead
Mejora continua & SPI-09 Acciones cerradas a tiempo; SPI-10 Auditorías ejecutadas & Muestra si el SMS corrige lo que encuentra y si se verifica cada año. \\ \hline
Cultura de seguridad & SPI-01 Reportes voluntarios; SPI-08 Personal con instrucción vigente & Un buen número de reportes y personal capacitado indican confianza y conocimiento del sistema. \\ \hline
Cumplimiento normativo & SPI-03 Eventos de reporte obligatorio; SPI-04 Accidentes e incidentes graves; SPI-07 Vuelos con FRAT completo & Mide que se cumpla el RAC 100 en cada vuelo y que no ocurran eventos graves. \\ \hline
Recursos adecuados & SPI-08 Personal con instrucción vigente; SPI-11 Discrepancias técnicas repetidas & Muestra si hay personal formado y equipos en condición para volar. \\ \hline
Operación segura en VLOS, vuelo nocturno y zona urbana & SPI-02 Errores del personal operativo; SPI-05 Salidas del área autorizada; SPI-06 Pérdidas de enlace C2 & Mide los riesgos propios de la operación certificada de Skyfilms. \\ \hline
\end{longtable}

\textbf{Programa Estatal de Seguridad Operacional (SSP):} los objetivos de Skyfilms apoyan los objetivos del Estado: prevenir accidentes, reportar los eventos exigidos y cumplir el RAC 100 y el RAC 219. Si la Aerocivil define indicadores o metas propios para explotadores UAS, Skyfilms los adopta y actualiza este anexo y el Anexo 9.

% =====================================================================
\clearpage
\section*{Anexo 2 --- Contactos de Emergencia}
\addcontentsline{toc}{section}{Anexo 2 --- Contactos de Emergencia}
La sede de Skyfilms está en Ibagué (Tolima). El Gerente SMS mantiene este directorio actualizado y lo revisa cada vez que cambia un dato y, como mínimo, cada seis meses. Antes de operar en otra ciudad o región, el Gerente SMS agrega los contactos locales a este anexo y los entrega al piloto al mando junto con la misión (Manual de Operaciones).

\begin{longtable}{|>{\RaggedRight\arraybackslash}p{5.2cm}|>{\RaggedRight\arraybackslash}p{4.6cm}|>{\RaggedRight\arraybackslash}p{5.8cm}|}
\hline
\rowcolor[gray]{0.9}\textbf{Contacto} & \textbf{Teléfono o canal} & \textbf{Cuándo se llama} \\ \hline
\endhead
Línea única de emergencias & 123 & Lesiones, incendio, accidente con personas o bienes de terceros \\ \hline
Bomberos & 119 & Incendio de batería o de equipo \\ \hline
Cruz Roja Colombiana & 132 & Apoyo en atención de personas lesionadas \\ \hline
Defensa Civil Colombiana & 144 & Apoyo en búsqueda de una aeronave extraviada o rescate \\ \hline
Aerocivil: canales oficiales de emergencia, publicados en la página web de la Aerocivil & Se consultan en la página oficial antes de cada operación & Accidente, incidente grave o aeronave fuera de control (ver el procedimiento de comunicación con la Aerocivil, capítulo 2) \\ \hline
Gerente General & [teléfono por diligenciar] & Activación del Plan de Respuesta a Emergencias \\ \hline
Gerente SMS & [teléfono por diligenciar] & Todo evento, para su reporte e investigación \\ \hline
Jefe de Pilotos UAS & [teléfono por diligenciar] & Evento durante una operación de vuelo \\ \hline
Técnico UAS (Responsable de Mantenimiento) & [teléfono por diligenciar] & Falla técnica, batería dañada o aeronave inoperativa \\ \hline
Centro médico más cercano al sitio de la misión & Lo define el piloto al mando en la planeación de la misión & Lesiones de personas \\ \hline
\end{longtable}

\textbf{Pruebas del directorio:} cada seis meses el Gerente SMS llama a cada número de la lista interna para comprobar que sigue vigente y deja constancia en el acta del SAG.

% =====================================================================
\clearpage
\section*{Anexo 3 --- Procedimiento de Control de Documentos y Registros}
\addcontentsline{toc}{section}{Anexo 3 --- Control de Documentos y Registros}

\textbf{Responsable único:} el Gerente SMS controla los documentos y registros del SMS. El Técnico UAS controla los del mantenimiento (MCM) y el Jefe de Pilotos los de operación (MO), y todos usan este mismo procedimiento.

\textbf{Lugar y conservación.} Los documentos vigentes están en el repositorio digital de Skyfilms (carpeta \textit{SMS}), con copia de respaldo mensual en un segundo medio, y en copia impresa controlada en la oficina de Skyfilms en Ibagué. Los registros se conservan \textbf{cinco (5) años}. Pasado ese plazo, el Gerente SMS los destruye o los archiva de forma permanente y lo anota en el listado maestro.

\textbf{Códigos.} Manuales: SMS-SKY-001, MO y MCM-SKY-001. Formatos: F-001 a F-022. Cada documento lleva código, versión y fecha.

\textbf{Matriz de acceso.} Las letras significan: \textbf{E} = elabora, \textbf{R} = revisa, \textbf{A} = aprueba, \textbf{C} = consulta, \textbf{S} = custodia (guarda el original).

\begin{longtable}{|>{\RaggedRight\arraybackslash}p{5.6cm}|c|c|c|c|c|}
\hline
\rowcolor[gray]{0.9}\textbf{Documento} & \textbf{Gerente General} & \textbf{Gerente SMS} & \textbf{Jefe de Pilotos} & \textbf{Técnico UAS} & \textbf{Pilotos} \\ \hline
\endhead
Manual SMS & A & E, S & R & C & C \\ \hline
Manual de Operaciones & A & R & E, S & C & C \\ \hline
Manual de Control de Mantenimiento & A & R & R & E, S & C \\ \hline
Reportes de peligros y eventos (F-013 y otros) & C & E, S & C & C & E \\ \hline
Registro de peligros y matrices de riesgo & A & E, S & R & C & C \\ \hline
Actas del SRB y del SAG & A (SRB) & E, S & C & C & -- \\ \hline
Registros de instrucción (F-016, F-017) & C & R & E, S & C & C \\ \hline
Registros de mantenimiento (F-018 a F-022) & C & R & C & E, S & C \\ \hline
\end{longtable}

\textbf{Pasos para cambiar un documento.} (1) El Gerente SMS o el dueño del documento elabora el cambio. (2) Lo revisan las personas con R. (3) El Gerente General lo aprueba. (4) El Gerente SMS cambia el número de versión y la fecha, y actualiza el listado maestro. (5) Distribuye la nueva versión y retira las obsoletas (acuse F-015). (6) Verifica en la siguiente revisión que no queden copias antiguas en uso.

\textbf{Documentos obsoletos.} Se marcan con la palabra OBSOLETO, se retiran de uso y se guarda una sola copia durante cinco años en una carpeta separada.

\textbf{Listado maestro.} El Gerente SMS mantiene una tabla con: código, título, versión, fecha de aprobación, quién aprobó, vigencia, lugar de custodia y tiempo de retención de cada documento y formato de Skyfilms.

% =====================================================================
\clearpage
\section*{Anexo 4 --- Listado de Reportes Obligatorios (MOR)}
\addcontentsline{toc}{section}{Anexo 4 --- Listado de Reportes Obligatorios (MOR)}
Se informa a la Aerocivil, por el Gerente General o el Gerente SMS, dentro de las 48 horas siguientes al evento (RAC 100.240; ver el procedimiento de reporte externo del Manual de Operaciones) cualquiera de estos eventos:

\begin{enumerate}
    \item Accidente: lesión grave o muerte de una persona, o daño importante a bienes de terceros o a la aeronave.
    \item Incidente grave: evento que casi termina en accidente.
    \item Pérdida de control de la aeronave o aeronave que se aleja sin control (fly-away).
    \item Pérdida del enlace de mando y control (C2) que lleve a un retorno automático o a perder la aeronave de vista.
    \item Salida del área o de la altura autorizadas.
    \item Cuasi colisión o acercamiento peligroso con una aeronave tripulada u otra aeronave.
    \item Colisión con personas, vehículos, estructuras, fauna u otros objetos.
    \item Incendio, humo o calentamiento anormal de una batería o de un equipo, o caída de la aeronave por ese motivo.
    \item Falla estructural o de un sistema crítico durante el vuelo (motor, hélice, controlador de vuelo).
    \item Interferencia que afecte la operación (señal, GPS u otra).
    \item Operación sin los permisos, autorizaciones o seguros exigidos, o fuera de las condiciones de las OpSpecs.
\end{enumerate}

Todo evento de esta lista se reporta en el F-013, se investiga formalmente (ver el capítulo 3) y lo revisa el SRB. Los eventos que no están en la lista se reportan como voluntarios.

% =====================================================================
\clearpage
\section*{Anexo 5 --- Guía de Uso del Formato de Notificación de Peligros}
\addcontentsline{toc}{section}{Anexo 5 --- Guía del Formato de Notificación de Peligros}
Se usa el F-013 (Reporte de Evento o Peligro). Lo puede diligenciar cualquier persona, también un tercero o un cliente, en papel, por correo o por mensaje al Gerente SMS.

\textbf{Qué escribir:} fecha y hora; lugar; aeronave; qué pasó o qué peligro vio; quién estaba; qué hizo en el momento; y qué sugiere. El nombre es opcional en reportes voluntarios.

\textbf{Plazos de gestión.}
\begin{longtable}{|>{\RaggedRight\arraybackslash}p{4.2cm}|>{\RaggedRight\arraybackslash}p{3.3cm}|>{\RaggedRight\arraybackslash}p{3.3cm}|>{\RaggedRight\arraybackslash}p{4.8cm}|}
\hline
\rowcolor[gray]{0.9}\textbf{Tipo de reporte} & \textbf{Acuse al reportante} & \textbf{Análisis inicial} & \textbf{Cierre y retroalimentación} \\ \hline
\endhead
Obligatorio (MOR) & 24 horas & 72 horas & Investigación en 30 días; retroalimentación al terminar \\ \hline
Voluntario con riesgo en zona roja o amarilla & 48 horas & 5 días hábiles & 20 días hábiles \\ \hline
Voluntario con riesgo verde & 5 días hábiles & 10 días hábiles & 30 días hábiles \\ \hline
\end{longtable}

\textbf{Protección de la información.}
\begin{itemize}
    \item Solo el Gerente SMS conoce el nombre del reportante. En informes, actas y boletines el nombre se quita (desidentificación): no se escriben nombres, matrículas ni lugares que permitan reconocerlo.
    \item El reporte no se usa para sancionar, salvo las excepciones de la Política (violación deliberada, negligencia grave, falsificación de registros u operar bajo alcohol o sustancias).
    \item Solo se entrega información de un reporte a un tercero cuando lo exija la Aerocivil u otra autoridad, o cuando haya un riesgo inminente para la vida de las personas. Aun así, se entrega la mínima información necesaria.
    \item La base de datos de reportes está en el repositorio digital con acceso solo del Gerente SMS y el Gerente General.
\end{itemize}

\textbf{Retroalimentación.} El Gerente SMS informa al reportante qué se encontró y qué se decidió, en persona, por llamada o por escrito, y lo comparte sin nombres en el boletín de seguridad del trimestre.

% =====================================================================
\clearpage
\section*{Anexo 6 --- Guía de Evaluación de Riesgos de Vuelo (FRAT)}
\addcontentsline{toc}{section}{Anexo 6 --- Guía FRAT}
El piloto al mando completa el FRAT antes de cada jornada de vuelo y lo entrega al Jefe de Pilotos. Se marca cada factor que aplique y se suman los puntos.

\begin{longtable}{|>{\RaggedRight\arraybackslash}p{4cm}|>{\RaggedRight\arraybackslash}p{9cm}|c|}
\hline
\rowcolor[gray]{0.9}\textbf{Grupo} & \textbf{Factor} & \textbf{Puntos} \\ \hline
\endhead
Personas & Menos de 7 horas de sueño, cansancio, estrés o malestar & 3 \\ \hline
Personas & Menos de 3 vuelos en los últimos 30 días con ese equipo & 2 \\ \hline
Personas & No hay observador o este no está descansado & 3 \\ \hline
Aeronave & Batería con más de 100 ciclos o con una discrepancia pendiente & 2 \\ \hline
Aeronave & Mantenimiento o actualización de firmware reciente sin prueba de vuelo & 3 \\ \hline
Aeronave & Un sistema de fábrica con falla (luces, sensores) & 5 \\ \hline
Entorno & Viento fuerte, lluvia, niebla o temperatura fuera de los límites del fabricante & 3 \\ \hline
Entorno & Vuelo nocturno & 2 \\ \hline
Entorno & Zona urbana, con personas o vehículos cerca & 3 \\ \hline
Entorno & Aeródromo, helipuerto o tráfico de aeronaves tripuladas cercano & 3 \\ \hline
Misión & Sitio nuevo o sin reconocimiento previo & 2 \\ \hline
Misión & Presión del cliente o poco tiempo para preparar el vuelo & 3 \\ \hline
\end{longtable}

\begin{tabularx}{\textwidth}{|l|X|}
\hline
\rowcolor[gray]{0.9}\textbf{Total} & \textbf{Decisión} \\ \hline
0 a 8 (verde) & El piloto al mando puede volar. \\ \hline
9 a 16 (amarillo) & Vuela solo con una medida de mitigación anotada y la autorización del Jefe de Pilotos. \\ \hline
17 o más (rojo) & No se vuela hasta reducir el puntaje; si no se puede, decide el Gerente SMS y se informa al Gerente General. \\ \hline
\end{tabularx}

\medskip
Un solo factor de 5 puntos (sistema de fábrica con falla) impide el vuelo, sin importar el total. Los FRAT se resumen cada trimestre para el SPI-07.

% =====================================================================
\clearpage
\section*{Anexo 7 --- Registro de Peligros de Seguridad Operacional}
\addcontentsline{toc}{section}{Anexo 7 --- Registro de Peligros}
Es el listado vivo de peligros de Skyfilms. El Gerente SMS lo actualiza con cada reporte, hallazgo o cambio y lo presenta al SRB. Aquí está la lista inicial para la operación de Skyfilms. La probabilidad es 1 a 5 y la gravedad A a E (ver la matriz del capítulo 3).

\begin{longtable}{|c|>{\RaggedRight\arraybackslash}p{3.4cm}|>{\RaggedRight\arraybackslash}p{3.1cm}|c|>{\RaggedRight\arraybackslash}p{4.5cm}|c|}
\hline
\rowcolor[gray]{0.9}\textbf{N.º} & \textbf{Peligro} & \textbf{Consecuencia} & \textbf{Riesgo inicial} & \textbf{Controles} & \textbf{Riesgo residual} \\ \hline
\endhead
1 & Pérdida del enlace C2 & Pérdida de la aeronave o caída & 3C & Prueba de enlace antes de volar; calidad de enlace mínima; retorno programado; observador & 2C \\ \hline
2 & Falla o incendio de batería & Incendio, lesiones, pérdida del equipo & 3B & Revisión y ciclos de baterías (MCM); bolsa ignífuga; extintor ABC; carga vigilada & 2B \\ \hline
3 & Colisión con personas o bienes de terceros en zona urbana & Lesiones o daños & 3A & Área despejada; no se sobrevuelan personas ajenas; altura y distancia mínimas; observador & 2A \\ \hline
4 & Vuelo nocturno: pérdida de orientación u obstáculos no vistos & Colisión o pérdida de la aeronave & 3C & Reconocimiento diurno previo; luces; límites de altura; observador & 2C \\ \hline
5 & FPV: campo de visión reducido del piloto & Colisión o desorientación & 3C & Observador siempre; distancia máxima de unos 250 m; reconocimiento del sitio & 2C \\ \hline
6 & Cercanía de aeronaves tripuladas & Cuasi colisión o colisión & 2A & Consulta de NOTAM y espacio aéreo; autorizaciones; suspensión ante cualquier aeronave & 1A \\ \hline
7 & Fatiga o error del piloto o del observador & Decisiones equivocadas & 3C & Descanso mínimo; FRAT; verificación cruzada; límites de turno & 2C \\ \hline
8 & Mantenimiento mal hecho o liberación sin pruebas & Falla en vuelo & 3B & Criterios de aceptación y rechazo; F-022; prueba funcional; revisión cruzada & 2B \\ \hline
9 & Mal clima inesperado (viento, lluvia) & Pérdida de control & 3C & Pronóstico antes del vuelo; límites del fabricante; suspensión ante cambios & 2C \\ \hline
10 & Cambio de normas, equipos o personal sin evaluar el riesgo & Incumplimiento o nuevos peligros & 3D & Procedimiento de gestión del cambio (Anexo 10) & 2D \\ \hline
\end{longtable}

\textit{Estos valores iniciales los ajusta el SAG con los reportes reales.}

% =====================================================================
\clearpage
\section*{Anexo 8 --- Matriz de Evaluación de Riesgos (Método Reactivo)}
\addcontentsline{toc}{section}{Anexo 8 --- Matriz de Evaluación de Riesgos (Reactivo)}
Se diligencia por cada evento reportado o investigado. Las matrices de probabilidad, gravedad y tolerabilidad están en el capítulo 3.

\begin{tabularx}{\textwidth}{|l|X|}
\hline
\rowcolor[gray]{0.9}\textbf{Campo} & \textbf{Qué se escribe} \\ \hline
Evento y fecha & Código del F-013 y fecha \\ \hline
Peligro y contexto & Qué es el peligro y en qué condición se dio \\ \hline
Causas y consecuencias & Causas, factores contribuyentes y qué pudo pasar \\ \hline
Supuestos o datos que faltan & Se escribe qué se supone y por qué; si falta información, se usa el peor caso creíble \\ \hline
Riesgo inicial & Probabilidad (1--5) y gravedad (A--E), por ejemplo 3C; región (roja, amarilla o verde) \\ \hline
Controles y responsable & Qué se hará, quién y para cuándo \\ \hline
Riesgo residual & Nuevo índice después de aplicar los controles \\ \hline
Aceptó el riesgo & Nombre, cargo y fecha (según el nivel de autoridad del capítulo 3) \\ \hline
Verificación de eficacia & Fecha, método y resultado (ver el capítulo 3) \\ \hline
\end{tabularx}

% =====================================================================
\clearpage
\section*{Anexo 9 --- Indicadores de Rendimiento de Seguridad Operacional (SPI)}
\addcontentsline{toc}{section}{Anexo 9 --- Indicadores de Rendimiento (SPI)}
Hay once indicadores. Se calculan cada trimestre, los revisa el SAG, y el SRB los analiza al menos una vez al año. Una desviación del nivel de alerta lleva a una acción correctiva en el Plan de Mejora Continua (Anexo 11).

\begin{longtable}{|c|>{\RaggedRight\arraybackslash}p{3.1cm}|>{\RaggedRight\arraybackslash}p{4.9cm}|>{\RaggedRight\arraybackslash}p{2.2cm}|>{\RaggedRight\arraybackslash}p{2.4cm}|}
\hline
\rowcolor[gray]{0.9}\textbf{N.º} & \textbf{Indicador} & \textbf{Fórmula (por trimestre)} & \textbf{Meta} & \textbf{Alerta} \\ \hline
\endhead
SPI-01 & Reportes voluntarios & Reportes voluntarios recibidos / vuelos realizados x 100 & Al menos 5 por cada 100 vuelos & Menos de 2 \\ \hline
SPI-02 & Errores del personal operativo & Errores reportados de pilotos y observadores / vuelos realizados x 100 & 5 o menos por cada 100 vuelos & Más de 10 \\ \hline
SPI-03 & Eventos de reporte obligatorio & Eventos del Anexo 4 / vuelos realizados x 100 & 0 & 1 o más \\ \hline
SPI-04 & Accidentes e incidentes graves & Número de accidentes e incidentes graves & 0 & 1 o más \\ \hline
SPI-05 & Salidas del área o altura autorizadas & Vuelos con salida del área o de la altura autorizada / vuelos realizados x 100 & 0 & 1 o más \\ \hline
SPI-06 & Pérdidas de enlace C2 & Pérdidas de enlace / vuelos realizados x 100 & 1 o menos por cada 100 vuelos & Más de 3 \\ \hline
SPI-07 & Vuelos con FRAT completo & Vuelos con FRAT completo / vuelos realizados x 100 & 100\,\% & Menos de 95\,\% \\ \hline
SPI-08 & Personal con instrucción vigente & Personas con instrucción SMS y recurrente vigente / personal total x 100 & 100\,\% & Menos de 90\,\% \\ \hline
SPI-09 & Acciones cerradas a tiempo & Acciones cerradas en su plazo / acciones con plazo vencido o cerrado x 100 & 90\,\% o más & Menos de 70\,\% \\ \hline
SPI-10 & Auditorías ejecutadas & Auditorías hechas / auditorías programadas x 100 & 100\,\% & Menos de 100\,\% \\ \hline
SPI-11 & Discrepancias técnicas repetidas & Discrepancias del mismo componente en la misma aeronave en el trimestre & 0 & 2 o más \\ \hline
\end{longtable}

\textbf{Periodo y denominador.} El periodo es un trimestre calendario. Si en un trimestre no hubo vuelos, los indicadores por 100 vuelos no se calculan y se anota ``sin operación''. Las metas de este anexo son las iniciales de Skyfilms y se ajustan con los datos de los primeros doce meses y con lo que defina la Aerocivil.

% =====================================================================
\clearpage
\section*{Anexo 10 --- Procedimiento de Gestión del Cambio}
\addcontentsline{toc}{section}{Anexo 10 --- Gestión del Cambio}
\textbf{Cuándo se aplica (desencadenantes).} Cuando cambie: (a) el tipo de operación, el sitio, la ciudad o el cliente; (b) una aeronave, equipo, software o firmware; (c) una persona en un cargo clave o personal nuevo; (d) un procedimiento o manual; (e) una norma del RAC 100, RAC 219 u otra aplicable; (f) un proveedor de mantenimiento; o (g) cuando un evento o auditoría lo exija.

\textbf{Pasos, responsables y plazos.}
\begin{longtable}{|c|>{\RaggedRight\arraybackslash}p{7.8cm}|>{\RaggedRight\arraybackslash}p{3.2cm}|>{\RaggedRight\arraybackslash}p{2.4cm}|}
\hline
\rowcolor[gray]{0.9}\textbf{Paso} & \textbf{Qué se hace} & \textbf{Responsable} & \textbf{Plazo} \\ \hline
\endhead
1 & Quien propone el cambio lo describe por escrito (qué, por qué, cuándo) y lo entrega al Gerente SMS & Quien lo propone & Antes de aplicarlo \\ \hline
2 & Se identifican las personas y áreas afectadas (pilotos, mantenimiento, clientes, proveedores) y se les informa & Gerente SMS & 5 días hábiles \\ \hline
3 & Se evalúan los peligros y riesgos del cambio, los transitorios (durante el cambio) y los acumulados con otros cambios, y se revisan las evaluaciones ya hechas & SAG & 10 días hábiles \\ \hline
4 & Se definen los controles y se aprueba el cambio antes de aplicarlo. Los cambios de riesgo alto los aprueba el Gerente General; los demás, el Gerente SMS & Gerente SMS / Gerente General & Antes de aplicarlo \\ \hline
5 & Se aplica el cambio, con capacitación o comunicación si hace falta, y se actualizan los manuales y las OpSpecs. Si el cambio afecta lo certificado, se informa a la Aerocivil antes de operar & Gerente SMS & Según el cambio \\ \hline
6 & Se hace seguimiento: revisión a los 30 días y a los 90 días, y se registra el resultado & Gerente SMS & 30 y 90 días \\ \hline
\end{longtable}
Registro: la tabla de cambios del Plan de Mejora Continua (Anexo 11) y el acta del SAG.

% =====================================================================
\clearpage
\section*{Anexo 11 --- Plan de Mejora Continua}
\addcontentsline{toc}{section}{Anexo 11 --- Plan de Mejora Continua}
Plantilla para registrar hallazgos de auditorías, eventos, SPI en alerta y cambios. El Gerente SMS la revisa cada mes y la presenta al SAG.

\begin{longtable}{|>{\RaggedRight\arraybackslash}p{0.8cm}|>{\RaggedRight\arraybackslash}p{2.4cm}|>{\RaggedRight\arraybackslash}p{2.8cm}|>{\RaggedRight\arraybackslash}p{2.8cm}|>{\RaggedRight\arraybackslash}p{1.9cm}|>{\RaggedRight\arraybackslash}p{1.9cm}|>{\RaggedRight\arraybackslash}p{1.9cm}|}
\hline
\rowcolor[gray]{0.9}\textbf{N.º} & \textbf{Origen y fecha} & \textbf{Hallazgo o peligro} & \textbf{Acción} & \textbf{Responsable} & \textbf{Plazo} & \textbf{Estado y eficacia} \\ \hline
\endhead
 & & & & & & \\[1.2cm] \hline
 & & & & & & \\[1.2cm] \hline
 & & & & & & \\[1.2cm] \hline
\end{longtable}

\textbf{Verificación de eficacia.} Cada acción tiene un método (revisar registros, repetir la inspección, mirar el SPI), una fecha de revisión y un criterio de éxito. El Gerente SMS no cierra una acción hasta comprobarlo y anota el resultado en esta tabla.

% =====================================================================
\clearpage
\section*{Anexo 12 --- Programa de Instrucción SMS}
\addcontentsline{toc}{section}{Anexo 12 --- Programa de Instrucción SMS}
\begin{longtable}{|>{\RaggedRight\arraybackslash}p{3cm}|>{\RaggedRight\arraybackslash}p{2.8cm}|>{\RaggedRight\arraybackslash}p{4.2cm}|>{\RaggedRight\arraybackslash}p{1.9cm}|>{\RaggedRight\arraybackslash}p{3.1cm}|}
\hline
\rowcolor[gray]{0.9}\textbf{Curso} & \textbf{Para quién} & \textbf{Contenido y objetivo} & \textbf{Duración y modalidad} & \textbf{Cuándo y evaluación} \\ \hline
\endhead
Inducción SMS & Todo el personal nuevo & Política, objetivos, estructura del SMS, cómo reportar, cultura justa, funciones de cada cargo & 3 h, presencial & Antes de operar; examen escrito con 80\,\% mínimo \\ \hline
Reentrenamiento SMS & Todo el personal & Repaso del SMS, eventos y lecciones aprendidas, cambios del Manual & 2 h, presencial o virtual & Cada 4 meses; examen corto \\ \hline
Gestión del riesgo y FRAT & Pilotos, observadores, Jefe de Pilotos & Peligros, matriz de riesgo, uso del FRAT, decisiones de vuelo & 2 h, práctico & Al iniciar y al repasar; ejercicio con casos \\ \hline
Factores humanos y habilidades no técnicas & Todo el personal operativo y de mantenimiento & Fatiga, error humano, comunicación, trabajo en equipo, toma de decisiones, cultura justa & 2 h, presencial & Cada 4 meses; casos prácticos \\ \hline
Emergencias y simulacro & Todo el personal & Plan de respuesta a emergencias, baterías, pérdida de control, llamadas a la Aerocivil & 2 h, práctico & Simulacro cada 6 meses; informe con resultados \\ \hline
Personal clave del SMS & Gerente SMS y Jefe de Pilotos & Investigación de eventos, auditoría, análisis de SPI, gestión del cambio & 8 h, curso externo o con instructor calificado & Al designarlos y cada 24 meses; certificado o constancia \\ \hline
\end{longtable}

\textbf{Instructores.} Dicta los cursos el Gerente SMS, o el Jefe de Pilotos en los de operación, con formación demostrable en SMS o aviación. Para los cursos del personal clave se puede acudir a un instructor externo.

\textbf{Registros.} Cada curso se registra en el F-016 (asistencia) y el F-017 (evaluación y cierre), con fecha, tema, instructor y resultado. Se archivan cinco años.

\textbf{Si el resultado no alcanza.} La persona repite el tema con apoyo del instructor en un plazo de 15 días y no cumple funciones clave hasta aprobar.

\textbf{Eficacia del programa.} El Gerente SMS compara cada año los resultados de las evaluaciones, los SPI-02 y SPI-08 y los reportes, y ajusta los contenidos.

% =====================================================================
\clearpage
\section*{Documentación de Referencia}
\addcontentsline{toc}{section}{Documentación de Referencia}
\begin{itemize}
    \item Doc 9859 de la OACI --- Manual de Gestión de la Seguridad Operacional.
    \item RAC 219 --- Gestión de la Seguridad Operacional.
    \item RAC 100 --- Operación de Sistemas de Aeronaves No Tripuladas (UAS).
    \item 100-MAUT-5.0-22-017 CI --- Implementación del SMS en explotadores UAS.
    \item 219-MAUT-1.0-22-005 CI --- Definición de indicadores de rendimiento en materia de seguridad operacional.
    \item Circular GIVC-1.0-22-015 --- Respuesta ante emergencias.
\end{itemize}
\label{fin_anexos}
